\chapter{Orthogonal Polynomials}
\label{ch:13}
\section{Three-term recurrence in terms of $A$, $K$, and $\Lambda$}
\label{sec:13-1}
Differentiating the generating function gives
\[
k_{n+1}(x,N)=xk_n(x,N)-n(N-n+1)k_{n-1}(x,N).
\]
The spatial equation is
\[
\frac{N-x}{2}\bigl[k_n(x+2,N)-k_n(x,N)\bigr]
+\frac{N+x}{2}\bigl[k_n(x-2,N)-k_n(x,N)\bigr]
=-2n k_n(x,N).
\]
These formulas fix the normalization for later operator treatments.


\begin{remark}[Terminal polynomial and trivial integral zeros]
\label{rem:terminal-integral-zeros}
Extend the monic centered recurrence
\[
k_{n+1}(x,N)=xk_n(x,N)-n(N-n+1)k_{n-1}(x,N)
\]
through $n=N$. The resulting terminal polynomial is
\[
k_{N+1}(x,N)=\prod_{j=0}^{N}(x-(N-2j)).
\]
Indeed, at each spectral point $x=N-2j$, the generating function has degree $N$, so its coefficient of degree $N+1$ vanishes. The recurrence identifies this coefficient with $k_{N+1}/(N+1)!$, and monicity determines the factorization. Pairing opposite roots gives
\[
k_{N+1}(x,N)=
\begin{cases}
x\displaystyle\prod_{r=1}^{m}(x^2-4r^2),&N=2m,\\
\displaystyle\prod_{r=0}^{m}(x^2-(2r+1)^2),&N=2m+1.
\end{cases}
\]
Empty products are one. These are the trivial integral zeros: the entire finite spectrum $-N,-N+2,\ldots,N$, with spacing two. In ordinary-coefficient normalization the corresponding centered polynomial is $k_{N+1}(x,N)/(N+1)!$. The terminal polynomial vanishes on the support and has zero squared norm; it is not an additional positive-norm member of the orthogonal system.
\end{remark}

\subsection{Matrix form of the two recurrences}
\label{sec:matrix-recurrence-intertwining}
With $A_N,R_N,L_N,\Lambda_N$ defined in
Section~\ref{sec:finite-sl2-matrices}, and
$\Phi_N=(K_n(j;N))$ as in
Section~\ref{sec:matrix-Phi-definition}, the degree and spatial
recurrences take the paired form
\begin{equation}\label{eq:ency-intertwining}
A_N\Phi_N=\Phi_N\Lambda_N,\qquad
\Phi_NA_N=\Lambda_N\Phi_N.
\end{equation}
The coefficient proof and entrywise recurrences are given in
Section~\ref{sec:2-2}; the tensor interpretation is in
Section~\ref{sec:8-1}.
Terms outside the degree or column range are omitted.
Thus the columns of $\Phi_N$ are right eigenvectors of $A_N$,
and its rows are left eigenvectors, with the equally spaced spectrum
$N,N-2,\ldots,-N$. This does not make $A_N$ symmetric in the
standard basis. Explicit examples appear in
Appendix~\ref{app:matrix-tables}.

\section{Nonsymmetric case}
\label{sec:13-2}
To avoid silently changing the symmetric normalization, introduce a new family. Let $0<p<1$, $q=1-p$, and define
\begin{equation}\label{eq:biased-generating}
(1+pv)^{N-j}(1-qv)^j=\sum_{n=0}^N Q_n^{(p)}(j;N)v^n.
\end{equation}
If $J\sim\operatorname{Bin}(N,p)$, a single factor takes the value $1+pv$ with probability $q$ and $1-qv$ with probability $p$. Its linear coefficient is centered and has variance $pq$. Therefore
\[
\mathbb E[Q_n^{(p)}(J;N)Q_m^{(p)}(J;N)]
=\delta_{nm}\binom Nn(pq)^n.
\]
This is the two-point instance of Section~\ref{sec:5-3}. At $p=q=1/2$, $Q_n^{(1/2)}=2^{-n}K_n$; the new family is not numerically identical to $K_n$.
Writing $a=pN-j$, logarithmic differentiation of \eqref{eq:biased-generating} gives
\[
(n+1)Q_{n+1}^{(p)}=[a-(p-q)n]Q_n^{(p)}-pq(N-n+1)Q_{n-1}^{(p)}.
\]
For example, $Q_0^{(p)}=1$, $Q_1^{(p)}=a$, and $2Q_2^{(p)}=a^2-(p-q)a-pqN$.

The reversible birth--death generator on $0\le j\le N$ is
\[
(\mathcal L_pf)(j)=p(N-j)[f(j+1)-f(j)]+qj[f(j-1)-f(j)].
\]
Boundary terms with zero rates are omitted. Binomial detailed balance follows from
$\pi_p(j)p(N-j)=\pi_p(j+1)q(j+1)$, where $\pi_p(j)=\binom Njp^jq^{N-j}$. Acting on the generating function gives $\mathcal L_pG=-v\partial_vG$, hence
\[
\mathcal L_pQ_n^{(p)}=-nQ_n^{(p)}.
\]
This generator uses a one-coordinate refresh rate convention. At $p=1/2$ it is half the sign-flip generator whose eigenvalues are $-2n$ in Section~\ref{sec:13-1}. The factor is a time normalization, not a contradiction.

